% unconventional.tex -- Part "The frontier": three chapters on unconventional
% routes, to sit before the distance chapter (ch:distance).
% Labels created: ch:unconvmap, ch:tensor, ch:learning and the sec:um-*,
% sec:tn-*, sec:ln-*, tab:um-*, tab:tn-*, tab:ln-* labels below.

\chapter{A map of unconventional routes and their lower bounds}\label{ch:unconvmap}

The previous chapters measured index calculus as the literature defines it:
a factor base, a decomposition system, a solver and linear algebra, priced
against matched rho. The programme also kept a family of leads that change
one of those ingredients more radically: an extra oracle, a tensor network
in place of the decomposition system, a learned component in the solver, a
statistic that might show that a factor base is not random, a lift to
another curve model or another field. The lead atlas files them as the
\emph{exotic} family (\art{crypto-autoresearcher:knowledge/leads/ecdlp/exotic/}).
This chapter maps the family, with one table of verdicts and standing, and
then gives the lower bounds that most of the family runs into. The two
lines with the most work behind them have their own chapters: the
tensor-contraction route (Chapter~\ref{ch:tensor}) and the learning and
pseudorandomness lines (Chapter~\ref{ch:learning}).

Three conventions hold for all three chapters.
\begin{itemize}
\item \emph{Standing.} Two items in this family carry a proof-carrying
  claim: \texttt{PCC-20261104-1df881} (the source audit of implicit tensor
  entries) and \texttt{PCC-20261104-458dda} (the structured-predicate
  integration report), both of standing \emph{unverified}
  (\S\ref{sec:pt-tensor}). Every other figure is an unverified assertion of
  the producer named by its source. Several studies name their own
  validator or red-team report; that is recorded where it exists, but it is
  not a proof-carrying-claim standing and none of these studies has a
  \texttt{VALIDATION.json} gate row. No row of the index-calculus
  leaderboard comes from this family.
\item \emph{Costs.} No line in this family produced a counted cost of a
  complete discrete logarithm against matched rho. Where a study gives a
  ratio, it is a stage or structure statistic, and the pricing band of
  \S\ref{sec:band} is not applicable. Under the iterative success criterion
  none of the results is an improvement of the index-calculus baseline.
\item \emph{Scope.} Most measurements are at toy size: $p\le2^{14}$ for the
  structure instruments, $n\le17$ for the binary SAT studies, $p\le2039$ for
  the rank experiment. A toy-size null excludes only effects large enough
  to be seen at that size.
\end{itemize}

\section{The exotic family in the atlas}\label{sec:um-atlas}

On 6 November the exotic directory held 32 lead records. Their
dispositions, read from the YAML files on crypto-autoresearcher main, are:
\begin{itemize}
\item 26 \texttt{open\_next\_step};
\item 2 \texttt{open\_blocked} (advice leakage, the idea corpus);
\item 2 \texttt{exhausted\_at\_scale} (the crypto-size candidate screen and
  H-PSEUDO);
\item 1 \texttt{closed\_obstructed} (label-valued augmented oracles);
\item 1 \texttt{closed\_dominated} (the holant decomposition-count network).
\end{itemize}
A disposition records what the programme does next, not whether the route
works: most of the open leads have a named obstruction for the version
that was tried and a narrower question left. Table~\ref{tab:um-map} lists
(with Table~\ref{tab:um-map2}) every record with its best measured result as
the record states it and the chapter that discusses it. The table is a map, not a ranking. Where the
record's best result is ``none'', nothing was measured.

The family overlaps with work recorded elsewhere in this treatise. The
August oracle closures, the canonical-lift statistic and xedni lifting are
in Chapter~\ref{ch:summer}; the transport, monodromy, oracle and lifting
obstructions are in Chapter~\ref{ch:obstructions}; the $k$-sum tables are
in \S\ref{sec:pg-ksum}; the external tensor preprint is reviewed in
\S\ref{sec:pt-tensor}; the H-PSEUDO correction is in Chapter~\ref{ch:summer}
and Table~\ref{tab:corr-xref}. Corrections specific to this family are
collected in Table~\ref{tab:corr-exotic}.

\begin{table}[tbp]
\centering\footnotesize
\begin{tabularx}{\linewidth}{@{}>{\raggedright\arraybackslash}p{0.27\linewidth}>{\raggedright\arraybackslash}p{0.12\linewidth}>{\raggedright\arraybackslash}X>{\raggedright\arraybackslash}p{0.1\linewidth}@{}}
\toprule
Lead (\texttt{LEAD-EXO-}) & Disposition & Best result recorded; obstruction & Where \\
\midrule
\multicolumn{4}{@{}l}{\emph{Oracles and query bounds}}\\
\path{augmented-oracle-ggm} & closed, obstructed & label-valued and public-data oracles are Shoup-simulable; toy split saving $B/A=0.0557$ (preliminary) & \S\ref{sec:um-oracles} \\
\path{coordinate-revealing-oracles} & open & split off 4 Nov; ``the GGM gives no bound'' for $x$, $y$, full point, $x$-halving, field-valued jet and net, $C_t$, degree-$D$ tests & \S\ref{sec:um-oracles} \\
\path{low-degree-query-lower-bound} & open & H1 false as stated, H1$'$ proved; full rank in 39 of 39 cells, $p=251$--$2039$; no toy-size violation & \S\ref{sec:um-lowdeg} \\
\path{low-degree-algebraic-query-bound} & open & theorem programme; none & \S\ref{sec:um-lowdeg} \\
\path{ksum-conjecture-s5-oracle} & open & a 4-SUM oracle on $\{P:x(P)<p^{1/5}\}$ needs $n^{\lceil k/2\rceil}$ under the $k$-SUM conjecture; none measured & \S\ref{sec:um-generic} \\
\path{four-list-ktree-join} & open & $j=2$ four-tree at exponent $0.4167$ closed for character filters, toy $N\le65539$ & \S\ref{sec:um-generic} \\
\path{invertible-hint-subset} & open & none; Lange--Winterhof and Kiltz--Winterhof statements still to be read & \S\ref{sec:um-generic} \\
\path{lange-winterhof-interpolation} & open & none; theorem not yet recorded & \S\ref{sec:um-generic} \\
\path{elliptic-code-list-decoding} & open & none; counting bound $\binom BMp^{k-M-1}\ge1$, $M=\sqrt{Bk}$ & \S\ref{sec:um-generic} \\
\path{hnp-advice-leakage} & blocked & Cheon walk exponent $0.23$ against rho's $0.50$ only when $[\alpha^i]G$ is leaked & \S\ref{sec:um-generic} \\
\addlinespace
\multicolumn{4}{@{}l}{\emph{Tensor networks}}\\
\path{hashed-factor-base-tensor-ecdl} & open & none: 34 correctness tests; ECC2K-130 and ECCp-131 exceed the exact-leaf budget without a relation & Ch.~\ref{ch:tensor} \\
\path{holant-decomposition-count-network} & closed, dominated & digit-QTT rank in $x$-labelling at the ceiling (1024 at $p=1021$), as random relabellings; NULL & \S\ref{sec:ln-instruments} \\
\path{gag-holant-semaev-chain} & open & none & \S\ref{sec:ln-instruments} \\
\addlinespace
\multicolumn{4}{@{}l}{\emph{Learning}}\\
\path{ml-learned-models} & open & learned WDSat order 16{,}097 conflicts against 3{,}092 (confounded, \S\ref{sec:ln-components}) & \S\ref{sec:ln-ceiling}--\ref{sec:ln-components} \\
\addlinespace
\multicolumn{4}{@{}l}{\emph{Pseudorandomness and instruments}}\\
\path{dl-character-sum-pseudorandom} & exhausted & H-PSEUDO: $R$ consistent with 1 against a negation-closed null & Ch.~\ref{ch:summer} \\
\path{quadratic-phase-small-x-leakage} & open & NULL at $N\approx2^{12}$--$2^{14}$ and one curve at $2^{16}$ & \S\ref{sec:ln-instruments} \\
\path{coordinate-statistic-leakage} & open & census lowered from strong to preliminary (\texttt{CORR-20261026-c94668}) & \S\ref{sec:ln-instruments} \\
\path{weak-class-instrument-calibration} & open & no instrument detects anomalous or small-$k$ curves; $\zeta$-stable base detected & \S\ref{sec:ln-weak} \\
\path{methodology-instruments} & open & none; false-negative rate of the screens unmeasured & \S\ref{sec:ln-weak} \\
\bottomrule
\end{tabularx}
\caption[]{The exotic lead records on 6 November, part 1 of 2. Disposition and best result as each record states them; all figures producer-reported and unverified unless the text says otherwise. Source: \art{crypto-autoresearcher:knowledge/leads/ecdlp/exotic/}, one YAML file per row (\texttt{LEAD-EXO-}\emph{row}\texttt{.yaml}).}
\label{tab:um-map}
\end{table}

\begin{table}[tbp]
\centering\footnotesize
\begin{tabularx}{\linewidth}{@{}>{\raggedright\arraybackslash}p{0.27\linewidth}>{\raggedright\arraybackslash}p{0.12\linewidth}>{\raggedright\arraybackslash}X>{\raggedright\arraybackslash}p{0.1\linewidth}@{}}
\toprule
Lead (\texttt{LEAD-EXO-}) & Disposition & Best result recorded; obstruction & Where \\
\midrule
\multicolumn{4}{@{}l}{\emph{Lifts, transports, models}}\\
\path{global-lift-height} & open & xedni exponent 3 for one family; short Mordell--Weil relations at $p\le31$, no decomposition step & Ch.~\ref{ch:obstructions} \\
\path{padic-jet-canonical-lift} & open & digit identity exact 1360/1360 on hard lifts, whose cost is the lift & \S\ref{sec:ln-models} \\
\path{elliptic-net-eds} & open & charged exponent $1.97$--$2.80$ against BSGS $1.17$--$1.32$ & \S\ref{sec:ln-models} \\
\path{curve-model-representation} & open & $d_{\rm reg}=2$ for Weierstrass and Edwards at $m=2,3$ (toy) & \S\ref{sec:ln-models} \\
\path{split-jacobian-projection} & open & none; Hasse--Weil order inflation closure & \S\ref{sec:ln-models} \\
\path{miller-pairing-cocycle} & open & none; one unreviewed proposal & \S\ref{sec:ln-models} \\
\path{torsor-deck-period} & open & none; deck-orbit router no-go & \S\ref{sec:ln-models} \\
\path{lattes-xline-dynamics} & open & none; forward iterates are group operations & \S\ref{sec:ln-models} \\
\path{lossy-object-propagation} & open & none; orbit argument and rigidity theorem & \S\ref{sec:ln-models} \\
\addlinespace
\multicolumn{4}{@{}l}{\emph{Screens, corpora and out of scope}}\\
\path{cryptoscale-candidate-screen} & exhausted & thirteen batches, no admissible candidate & \S\ref{sec:um-routes} \\
\path{ideas-corpus-algebraic-transforms} & blocked & none; 15 active registry rows lack constructors & \S\ref{sec:um-routes} \\
\path{minors-guess-and-determine} & open & none; design reconstructed from abstracts & \S\ref{sec:um-routes} \\
\path{quantum-ecdlp-resources} & open & reversible secp256k1 addition, $-0.124\%$ on a sample island; not classical & \S\ref{sec:ln-scope} \\
\bottomrule
\end{tabularx}
\caption[]{The exotic lead records, part 2 of 2 (Table~\ref{tab:um-map} continued). Same source and standing.}
\label{tab:um-map2}
\end{table}

\section{Routes by which an encoding can be read}\label{sec:um-routes}

A sub-exponential algorithm for the ECDLP on a prime-order subgroup cannot
be generic. Shoup's bound puts every algorithm that uses only the group law
and equality at $\Omega(\sqrt N)$, and with $S$ bits of advice the
preprocessing bound is $ST^2\gtrsim N$~\cite{cgk2018}. A candidate must
therefore read the encoding: coordinates, the curve equation, the field.
The sub-exponential desk evaluation of 1 November
(\art{research/ecdlp_subexp_desk_20261101/README.md}) sorted the routes by
which an encoding can be read and matched sixteen proposed mechanisms
against the atlas with \texttt{lead\_atlas.py --match}; every one of the
sixteen landed on an existing lead. Its route table, as recorded at atlas
snapshot \texttt{cfc1dae2}, gives one named obstruction per route:
\begin{itemize}
\item an algebraic factor base of size $B$ over $\FF_p$ has description
  degree at least $B/3$, and mean yield is $\binom{B+m-1}m/N$ for every base
  of size $B$ (Chapter~\ref{ch:cover});
\item Coppersmith lattices reach at most $0.3293$ per $\Sem3$ block in three
  unknowns against a demand of $0.8889$ at $(m,w)=(4,2)$;
\item the digit and subset-sum form uses only group operations and
  equality, so it is generic;
\item enumerative, pair-table and meet-in-the-middle oracles have no
  sub-rho cell in the replicated cost model at any arity;
\item $p$-adic lifts, jets, canonical and Teichm\"uller lifts carry no
  additive image of order $\ell$ when $\gcd(n,p)=1$;
\item homomorphic images in a multiplicative group factor through the
  embedding degree;
\item index calculus in a genus-$g$ Jacobian over $\FF_p$ costs
  $p^{2-2/g}\ge p$ for $g\ge2$;
\item generic cost is an isogeny-class invariant, and automorphisms give
  $\sqrt{\abs{\mathrm{Aut}}}$.
\end{itemize}
The README now carries a banner: it is a historical snapshot of 154 leads,
a later reconciliation catalogued 160, and its ``xedni is closed''
statements are superseded by the conditional reading of
Chapter~\ref{ch:obstructions}. The route table is otherwise unchanged.

Three searches were broad screens rather than mechanisms. The crypto-size
candidate screen (\texttt{GOAL-CRYPTO-001}) ran thirteen batches that ended
\path{FINAL_NO_CANDIDATE} or \path{NO_ADMISSIBLE_NEXT_MECHANISM} over
its checked source set (\texttt{EV-CRYPTO-010}..\texttt{013}); its record
states the scale reached as ``zero-compute admission sheets; no crypto-size
run'' and asks for a sensitivity audit showing that the screen could have
admitted a known positive. The idea corpus has 15 active registry rows,
each blocked on a constructor or theorem that no one supplied. The minors
and guess-and-determine lead, which would extend the lit2026 W12
reproduction of \S\ref{sec:prime}, has a draft design reconstructed from
abstract pages only. None of the three produced a candidate to measure.

\section{Generic and conjecture-level obstructions}\label{sec:um-generic}

Five records in Table~\ref{tab:um-map} are bounded by a statement about
generic or structureless algorithms rather than by a measurement.
\begin{itemize}
\item \emph{$k$-SUM.} The prime-field tier that needs a 4-SUM oracle on
  $S_B=\{P:x(P)<p^{1/5}\}$ in time well below $\sqrt{\abs{S_B}}$ meets the
  $k$-SUM conjecture: structureless $k$-SUM needs $n^{\lceil k/2\rceil}$. The
  one known escape, clustered 3SUM, needs additive energy that the
  coincidence and character-sum measurements exclude. The record's next step
  is to write the reduction as a conditional obstruction, with the caveat
  that $k$-SUM has near-linear linear-decision-tree complexity, so no
  Shoup-style proof is available. The measured $k$-sum tables are in
  \S\ref{sec:pg-ksum}.
\item \emph{Four-list joins.} A Wagner tree on $E(\FF_p)$ needs a cheap
  filter $h$ with $h(P+Q)=f(h(P),h(Q))$ often enough. For character filters
  under group-law or quasigroup combining, the $j=2$ four-tree at exponent
  $0.4167$ is closed at toy $N\le65539$
  (\art{crypto-autoresearcher:analysis/o2-sum-compatible-filters/}); the
  log-domain FFT scans covered 507 families and 13{,}528 combinations, with a
  maximum of $1.10\times$ the null
  (\art{research/ic_prime_eth_leads_20261031/SYNTHESIS.md}, proposal 8). A
  planted log-mod-$M$ character is the Pohlig--Hellman leak, not a
  coordinate-computable filter.
\item \emph{Invertible hint subsets and interpolation degree.} A dense
  subset of $E(\FF_p)$ whose membership and logarithms are both computable
  from coordinates would beat the square-root bound by random
  self-reduction. The Lange--Winterhof and Kiltz--Winterhof degree bounds on
  polynomials interpolating the ECDL and ECDH maps are the expected closure,
  but both records say the theorem statements have not been read beyond
  their abstracts; whether they cover the statistics of
  Chapter~\ref{ch:learning} is open.
\item \emph{List decoding.} Using Guruswami--Sudan list decoding of elliptic
  codes as a decomposition oracle predicts hits only when
  $\binom BMp^{k-M-1}\ge1$ with $M=\sqrt{Bk}$, which forces $k$ of order $B$;
  nothing was run.
\item \emph{Advice.} The hidden-number conversion of a log-bit leak is
  deferred until some statistic has a measured nonzero advantage. The Cheon
  walk on leaked powers $[\alpha^i]G$ measured exponent $0.23$ against rho's
  $0.50$ from 24 bits up, but that is a different problem (the discrete
  logarithm with auxiliary inputs), not index calculus.
\end{itemize}

\section{Augmented oracles after the 4 November split}\label{sec:um-oracles}

Chapter~\ref{ch:obstructions} records the August result that jet,
endomorphism, incidence, halving and net oracles are simulable in the
generic group model. A red-team review of 4 November
(\art{crypto-autoresearcher:docs/reports/leads-round2-20261104/augmented-oracle-review.md})
showed that the record used ``oracle'' for three kinds of object, and the
owner split the lead into two the same day:
\begin{itemize}
\item \emph{label oracles} (input a label, output a label or a bit fixed by
  group relations): endomorphism image, label-valued halving, incidence
  check over a label factor base, label-valued net. These are Shoup-simulable
  with overhead $\Theta(\log N)$ or $O(m)$, and stay closed in
  \texttt{LEAD-EXO-augmented-oracle-ggm};
\item \emph{coordinate oracles} (input a label, output field data):
  $x$, $y$, full point, $x$-valued halving, field-valued jet and net, the
  threshold oracle $C_t$, degree-$D$ coordinate tests. These moved to
  \texttt{LEAD-EXO-coordinate-revealing-oracles}, whose record says ``the GGM
  gives no bound'';
\item \emph{public-data ``oracles''} with no instance input (factor-base
  membership of a field element): not augmentations at all, 0 queries.
\end{itemize}
The key step is Lemma~3 of the holant note
(\art{crypto-autoresearcher:docs/reports/holant-h1-ggm-lemma-20261103/README.md}):
on every curve the $x$-oracle cannot be simulated with $o(\sqrt{N'})$ group
queries, $N'$ the largest prime factor of $N$. Two $x$-queries give $x(G)$
and $x(Q)$, free local brute force gives $k$ up to sign, and $O(\log N)$
group queries fix the sign. The review's sanity check: if a full-point
oracle were simulable with $o(\sqrt N)$ overhead, Shoup would bound concrete
ECDLP by $\Omega(\sqrt N)$ on every curve, which Smart's attack on anomalous
curves contradicts. The same review found the endomorphism-image argument's
premise ``$\mathrm{End}_{\FF_p}(E)=\ZZ$'' false (Frobenius) and replaced it with
an eigenvalue argument with the same exponent, and it read the earlier jet
verdict ``$C=1$'' as the structured-GGM verdict withdrawn by
\texttt{CORR-20260727-003} (Table~\ref{tab:corr-exotic}).

The consequence for index calculus is the point of the split. Every
index-calculus algorithm reads coordinates, so it is a
coordinate-revealing algorithm, and the generic group model says nothing
about it. What the open lead asks is whether coordinate access gives a
\emph{time-bounded} advantage on prime-order prime-field curves, which is
the index-calculus question itself (\texttt{KN-OPEN-001}). The only
measured number on either side is a toy split saving: on 7--8-bit curves at
$m=3$, an $x$-membership-filtered Semaev meet in the middle has the same
yield as the unfiltered one in 40 of 40 cells and costs $0.0557$ of its field
operations, recorded as preliminary and as ``not a bound''.

\section{The low-degree query lower bound}\label{sec:um-lowdeg}

One way to keep a lower bound while admitting coordinate access is to
allow only low-degree tests: the solver may ask whether
$f(\text{coordinates of }X_1,\dots,X_t)=0$ for a polynomial $f$ of total
degree at most $D$, with $X_i=[a_i]P+[b_i]Q$. The lead's conjecture H1 was
success probability $O(T^2/N+TD/\sqrt N)$. A desk note of 4 November
(\art{crypto-autoresearcher:docs/reports/ecdlp-leads-desk-20261104/low-degree-desk.md},
claims tagged proved, cited or heuristic) settled the model.
\begin{itemize}
\item \emph{Single-point tests do not leak} (proved): a degree-$D$ test on
  one point is true for at most $2D$ ($3D$ with $y$) values of $k$.
\item \emph{H1 as stated is false} (proved, given a rank condition). The
  pulled-back function has degree up to
  $\Delta=\sum_i(2d_{x,i}+3d_{y,i})b_i^2$ on $E$, and large $b_i$ cost only
  $O(\log N)$ group operations. If the evaluation matrix of degree-$\le D$
  monomials on the $N$ points $v(R)$ has full rank, one test answers
  ``$k\in S$?'' for any $S$, and binary search finds $k$ with
  $\lceil\log_2N\rceil$ tests; for $D=3$ this needs $t\approx(6N)^{1/3}$
  points, so $T=O(N^{1/3}\log N)$.
\item \emph{The restatement H1$'$ is proved}: success is at most
  $\bigl(\binom{T_g+2}2+\sum_j(\Delta_j+t_j)+1\bigr)/N$, so sub-$\sqrt N$
  success needs $\sum_j\Delta_j\gtrsim N$. The bound is meaningful for
  low-degree statistics of points with small multipliers.
\item \emph{The Kohel--Shparlinski bound does not suffice} for $t\ge2$ with
  large multipliers: every bound in that paper scales with the degree of the
  pulled-back function, not with $D$. The note read the paper (ANTS-IV,
  2000) in full and records what it proves.
\end{itemize}

The rank condition was then tested exactly under a protocol frozen before
data (\art{crypto-autoresearcher:docs/reports/lowdeg-rank-20261104/README.md},
\texttt{LOWDEG-RANK-20261104}). With random multipliers, full rank held in
39 of 39 decisive cells at $p=251$, $509$, $1021$ and $2039$, once the
monomial count was about twice the number of candidates; the
random-point null was full rank and the shift control stayed within its
Riemann--Roch bound $3Dt+1$. Small multipliers $b_i=1..t$ failed in exactly
the three cells with $\sum b_i^2\approx0.34$--$0.41N$, as H1$'$ predicts. The
report's validator re-ranked 8 cells exactly, confirmed the verdict, and
corrected the protocol's own reading: at the tested points the attack spends
about 180--490 group operations, more than $\sqrt N$ (32--45), so \emph{no
toy-size violation of H1's query bound is shown}. By the validator's
accounting the attack would beat $\sqrt N/D$ only for $D\ge3$ and $N$ from
about $2^{28}$ ($D=5$) to $2^{48}$ ($D=3$), and never for $D=2$, which rests on
full rank persisting at $t\sim N^{1/D}$, untested.

The attack is not a cost. Building the test takes about $N^\omega$ local
field operations, and a test that answers ``$k\in S$?'' for $\abs S\approx N/2$
needs a polynomial with about $N/2$ monomials, so a plain logarithm table
dominates it. The next step on the lead is a bound that charges the test's
description length or local work, which, the desk note observes, would be a
complexity conjecture rather than a generic-model theorem. For index
calculus the lesson is the same as in \S\ref{sec:um-oracles}: once local work
on coordinates is free, query bounds stop constraining what a factor base
can be.

\section{The lit2026 W3 audit against generic bounds}\label{sec:um-w3}

The third bound is a check of the programme's own records against three
recent generic lower bounds (\art{research/lit2026_w3_ggm_audit_20261011/README.md},
workstream W3 of the lit2026 programme, 11 October; the other workstreams
are in \S\ref{sec:prime}). The three papers are:
\begin{itemize}
\item the structured generic-group model (SGGM) of Corrigan-Gibbs, Henzinger
  and Wu (ePrint 2026/384): an algorithm that exploits structure on a
  $\delta$ fraction of the group needs $\Omega(\min(\sqrt q,1/\delta))$ group
  operations;
\item the non-adaptive preprocessing bound of Dinur, Keller and Marmor
  (arXiv 2505.00894): success at most $3T^2/N+4\ln2\cdot ST/N$, so
  non-adaptive algorithms below $T=\sqrt N$ need $ST\gtrsim N$;
\item the multi-instance bound of Akshima, Kiltz, Nishaburi, Nooripoor and
  Wiedijk (ePrint 2026/2253), read in abstract only.
\end{itemize}
The study maps an $m$-term decomposition oracle over a factor base $F$ onto
the SGGM's partial operation $\star$, so that the constrained labels are $F$
and its $j$-sums, $j\le m$, and
$\delta\le c\sum_{j=1}^m\binom{\abs F+j-1}j/r$, with $c$ the cofactor (loose) or
1. That mapping is the study's derivation, not a theorem of the paper, and
the elliptic-curve section of 2026/384 could not be extracted from the copy
that was read (a stated gap).

\begin{table}[t]
\centering\small
\begin{tabularx}{\linewidth}{@{}Xrrrr@{}}
\toprule
Family & Rows & $\delta_{\rm upper}<1$ & min measured / bound & below a bound \\
\midrule
prime (Table B) & 72 & 48 & 38.9 & 0 \\
random binary (Table B) & 35 & 23 & 56.8 & 0 \\
Koblitz round 5 ($n\le41$) & 72 & 13 & 355 & 0 \\
Koblitz one target (Table A) & 6 & 1 & $2.84\cdot10^6$ & 0 \\
Koblitz 32-target batch (Table C) & 9 & 1 & $7.99\cdot10^4$ & 0 \\
Koblitz compact orbit $n37$--$n83$ (probes) & 14 & 0 & 1077 & 0 \\
\midrule
All & 208 & 86 & & 0 \\
\bottomrule
\end{tabularx}
\caption[]{The lit2026 W3 audit: 208 index-calculus and precompute
configurations against the SGGM bound. Producer-reported, from committed
records, isolation L0; the compact-orbit rows are in probes, not group
operations. Source: \art{research/lit2026_w3_ggm_audit_20261011/README.md},
Results \S1.}
\label{tab:um-w3}
\end{table}

Table~\ref{tab:um-w3} gives the result. No configuration is below any
bound, for $\min(\sqrt{r/A},1/\delta)$, for the explicit constant of the
paper's Theorem~3.2 at $\varepsilon=1/2$, or for the uniform-landing
$\delta$. The closest row,
\texttt{icv1-fp14-tm43-158b6914/mitm\_m2\_negfold\_walk\_balanced}, is
$38.9\times$ above its bound, and the balanced $m=2$ prime rows at 10--18
bits are $39$--$52\times$ above. The informative number is the other one:
\emph{122 of 208 rows have $\delta_{\rm upper}=1$}. Their factor base's
$m$-sums cover the subgroup, so the bound is 1 and certifies nothing; this
covers every compact-orbit rung and every Koblitz $m=3$ row with $c=h$. For
those configurations the whole cost is in \emph{finding} a decomposition,
which the SGGM treats as a free evaluation of $\star$. The SGGM therefore
splits the programme's configurations into sparse ones, which it bounds and
which all sit $39\times$ or more above the bound, and dense ones, which it
cannot see. Thirteen rows show relations per trial above the counting
ceiling, up to $1.2\times$; none is significant at the Bonferroni level
(smallest $p=0.0087$), so the study reads them as noise, not as an
accounting bug.

The other two papers place the table arms. Every non-adaptive table arm in
the repository sits on $ST\approx N/A$: the multiplier-closure arms at
$0.739$--$1.432$ (canonical entries times $T$ over $r/A$, 27 arms) and the
compact-orbit target extraction at $1.053$--$2.734$. Dinur--Keller--Marmor
say that is the floor for the class. Batch rho follows the
$\sqrt{kq}$ multi-instance baseline: fitted slopes $0.521$--$0.585$ with
every bootstrap interval covering the Kuhn--Struik finite-$k$ slope
$0.533$--$0.539$, on 744 plus 744 verified fresh logarithms. The audit also
found three record problems (an $n=61$ subgroup size recorded as
$\log_2r=47.5407$ instead of $47.2109$; excluded precompute defined
differently by rung; two frozen targets at $n71$ and $n73$ about
$0.0065$--$0.0075\times$ as expensive as the rank-stage mean, each a
$P\approx0.006$--$0.007$ event); they are reported, not edited.

\section{What the bounds leave}\label{sec:um-reading}

The three bounds agree. Generic and query models bound the label side of
an algorithm and are silent on the side index calculus works on: free
local work with coordinates. The SGGM is vacuous on the dense
configurations that carry the cost, the augmented-oracle split moves every
coordinate oracle outside Shoup's model, and the low-degree bound holds
only for small multipliers, because free local interpolation breaks it
otherwise. None of the three is evidence for or against index calculus
below rho; each says where a lower bound would have to charge local work
to say anything. The open leads in Table~\ref{tab:um-map} are, in effect,
attempts either to find such a charge or to find local work that a charge
would not cover. The next chapter is the attempt with the strongest claim
behind it.

% ---------------------------------------------------------------------------

\chapter{The tensor-contraction route}\label{ch:tensor}

\S\ref{sec:pt-tensor} reviews the external claim in brief: the preprint of
Akshima and coauthors, \emph{A Subexponential Time Algorithm for Elliptic
Curve Discrete Logarithm}, version~0.2 (Bochum/MPI, pinned at
\texttt{1ee44e56}), claims prime-field ECDLP in
$2^{O(b/(\log b)^{1/4})}$, $b=\lceil\log_2p\rceil$, by hashed factor bases and
exact Boolean tensor contraction. This chapter gives the full account of
what the programme built and learned on that route: the dependence on the
external theorem, the prototype and the native implementation, the
challenge attempts, the twelve prime-field leads and their refutations,
and the finite-size situation. The summary is one sentence. \emph{Nothing
on this route has a counted cost: no relation was found on ECC2K-130 or
ECCp-131, no exponent moved, and the route's sub-exponential claim is
conditional on an external result that is unverified.}

\section{What the claim rests on}\label{sec:tn-conditional}

The Bochum preprint's Theorem~4 is conditional on its Theorem~2, which
invokes Theorem~D.1 of the OpenAI preprint \emph{The Exponential Time
Hypothesis is False} (17 October 2026, revised 25 October; pinned at
\texttt{openai/math} \texttt{7ef5c0cc}). D.1 contracts an $s$-edge
bipartite Boolean tensor network, with degree at most $\lceil s^{9/10}\rceil$
and vertex tensors evaluated as black boxes, in
$(s+L)^{O(1)}2^{O(s/\sqrt{\log s})}$ time and polynomial space
(\art{research/ic_prime_eth_leads_20261031/LEADS.md} \S1).

\paragraph{The external result is unverified.} The programme did not
verify D.1, and as of the sources read here no one had. The lead note of 31
October records, from the public discussion it consulted, that no bug had
been reported, that independent experts were examining the proof, that the
Lean development in \texttt{openai/math} proves a weaker rate
($n(\log\log n)^2/\sqrt{\log n}$) and does not contain Theorem~D.1, and that
nobody had machine-checked that the formal statements match the paper; its
status line is ``claimed, under examination, no reported bug''. The
atlas's literature note on the preprint is tagged
\texttt{unverified-proof} and states its proof status as ``not checked here''
(\art{crypto-autoresearcher:knowledge/literature/KN-LIT-caa51d.md}). The two
in-session readings of D.1 and the D.4 recurrence described in
\S\ref{sec:pt-tensor} found no concrete gap; they are scoped readings, not a
certification. Every statement in this chapter about the asymptotic
behaviour of the route therefore inherits that condition, and a reader
should treat the route's sub-exponential bound as conditional on an
unverified theorem.

\paragraph{Where D.1 stops being inert.} D.1's exponential term depends on
the edge count only: vertices are black boxes, so a linear or low-rank
vertex gets no discount, and edges are bits, so a $p$-ary digit network
would need a $p$-adic analogue the paper does not supply
(\art{research/ic_prime_eth_leads_20261031/LEADS.md} \S1, facts 1--3). The
recursion's admission condition $u_0=\lambda/2^{17}\ge\lambda^{4/5}$ needs
$\lambda\ge2^{85}$, that is $s\ge2^{2^{85}}$ edges at depth 0 (derived by a
delegated audit; the Bochum paper quotes the weaker $2^{2^{17}}$). Below
that threshold the algorithm is brute force over the network. The native
implementation states the same fact in its own constants: the admission
scale $u=\lfloor\lceil\log_2s\rceil/131072\rfloor$ is zero for every
practically addressable network, so the code always selects the
manuscript's finite-size fallback (\art{docs/hashed-ecdl.md}).

\section{The prototype}\label{sec:tn-prototype}

The first build (\art{research/btn_ecdl_20261029/README.md}, 29 October) was
a Python prototype of the relation-finding step on the Koblitz curve
$y^2+xy=x^3+1$, with a plain bucket-elimination contractor standing in for
the claimed algorithm. Its findings, all local and seeded:
\begin{itemize}
\item over $\FF_{2^5}$ the network count of ordered pairs
  $A_1+A_2=R$ equals brute force (42 and 40 on two targets);
\item over $\FF_{2^7}$ the count is 112 against a brute-force 112, but plain
  elimination needs tables of $2^{34}$ entries (min-fill width 34, against
  23 at $n=5$), so the count conditions on the 7 bits of $x_1$ and takes
  336\,s; dropping the curve-membership predicate on $P_2$ as redundant
  gives 237 instead of 112, a recorded error;
\item a complete $m=2$ index calculus over $\FF_{2^5}$ ($r=11$, $N=16$) found
  17 relations and recovered the planted logarithm; one contraction took
  about 340\,s;
\item on the ECC2K-130 field the addition network with two
  curve-membership predicates has 105{,}062 variables, 104{,}801 predicates
  and 524{,}008 edges; a witness from reference arithmetic satisfies all
  predicates and a one-bit flip violates one. Factor-base predicates are
  omitted at that size. This is an encoding check, not an attack.
\end{itemize}
The lead note reads the widths as $2^{5n}$, five times the enumeration
exponent. The README lists what the prototype does not show: the multiplier
chain has $O(n^2)$ edges, the paper's $O(n)$ Fourier multiplier does not carry
over to $\FF_2[x]$ without a characteristic-2 variant, and the GPU
contractor was never run on a GPU.

\section{The native implementation and the challenge attempts}\label{sec:tn-native}

The Rust implementation of the complete algorithm, with the D.4 recursive
backend, is in \path{crates/hashed-ecdl-tensor} and mirrored in the
cryptanalysis repository (\art{docs/hashed-ecdl.md}). Its correctness
record is in \S\ref{sec:pt-tensor} (34 release tests, forced recursive and
direct controls, a tiny complete relation pipeline). Table~\ref{tab:tn-native}
lists the engineering studies that followed and what each did and did not
establish.

\begin{table}[t]
\centering\small
\begin{tabularx}{\linewidth}{@{}>{\raggedright\arraybackslash}p{0.26\linewidth}X@{}}
\toprule
Study & Recorded result \\
\midrule
ECCp-131 input validation, 29 Oct & public parameters as runner input; 27 regression tests pass; the challenge command exits with \texttt{resource limit: AKS multiplicative order} before graph or relation search \\
Exact primality and functional witness, 30 Oct & AKS kept above \texttt{u64}, Pocklington proof nodes verified; ECCp-131 validation at 5{,}589 charged diagnostic steps; 34 tests; matched performance evaluation blocked \\
Checked hash prime, 30 Oct & a user-reported 54-minute set-up hang traced to the AKS branch above \texttt{u64}, not to Fourier-network construction; attribution corrected \\
Primality fast paths, 30 Oct & thirteen strong witnesses below $3317044064679887385961981$, bounded Pocklington; original ECCp-131 input validates in 17{,}471 charged steps; 74 package tests; the challenge network constructs with 1{,}684{,}796 edges, 249{,}620 vertices, degree 1{,}032, 526 interface bits \\
CPU oracle panels, 30 Oct & two matched panels of 120 executions, byte-identical certificates; median $45.99\to3.90$\,ms ($p=65521$, $m=3$, $\tau=80$) and $80.54\to4.12$\,ms ($m=4$, $\tau=64$); L0 on one EKS VM; the oracle is enumeration with cached MITM tables, not contraction \\
Counts-only protocol, 31 Oct & v2 supersedes v1, withdrawn before data; isolation forced from L2 to L0 because rho solves under 1\,ms grade L1; no session has run \\
Entry-cost audit, 4 Nov & \texttt{PCC-20261104-1df881}, unverified: D.1 permits an exponential entry-work allowance; maps target-row budgets to entry traversal \\
\bottomrule
\end{tabularx}
\caption[]{Engineering on the native hashed-ECDL implementation. All
charged-step counts are uncalibrated admission diagnostics, not costs.
Producer-reported. Sources: \art{docs/hashed-ecdl-eccp131-validation.md};
\art{research/hashed_ecdl_speed_20261029/};
\art{research/hashed_ecdl_hash_prime_20261030/};
\art{research/hashed_ecdl_primality_fast_paths_20261030/};
\art{research/hashed_ecdl_cpu_20261030/};
\art{research/hashed_ecdl_counts_20261031/};
\art{research/hashed_ecdl_entry_cost_20261104/}.}
\label{tab:tn-native}
\end{table}

\paragraph{No relation on either challenge.} Two attempts at the published
challenges stopped at resource limits. The ECCp-131 network compiles to
1{,}564{,}506 edges, 654{,}523 variables and degree 1{,}032 in the Python
companion, which reaches its default $u_0=0$ exact-leaf budget without a
witness; the native command first stopped in prime certification, and after
the primality fixes constructs the network but has not produced a
relation. The ECC2K-130 network has 2{,}694{,}680 edges and also exceeds the
exact-leaf budget without a relation. The atlas records the state as
``ECC2K-130 and ECCp-131 attempts exceed their exact-leaf budget without a
relation; no runtime, $S$ or $S$/rho figure''
(\art{crypto-autoresearcher:knowledge/leads/ecdlp/exotic/LEAD-EXO-hashed-factor-base-tensor-ecdl.yaml};
\art{docs/hashed-ecdl.md}). The validation receipt states that the budget
failures ``do not prove nonexistence, a practical speedup, or a theorem
obstruction''. Both challenge scalars remain unresolved, and the
independent correctness and complexity audit of the backend is the open
problem \art{problems/hashed-ecdl-tensor-backend.md}.

\paragraph{The speed panels are not tensor measurements.} The CPU panels of
Table~\ref{tab:tn-native} are the one place on this route with matched
timings. Their $11.8$--$19.6\times$ median gain at $p=65521$ is the cached
meet-in-the-middle oracle against the original enumeration oracle, on
three public known-answer fixtures; the README states that the panels do not
establish a matched-rho advantage or ``the paper's conditional tensor
contraction hypothesis''. The larger-field feasibility ladder expired at its
scheduling deadline with no rows. The Rust port of the Go prototype and its
ecbench arm \texttt{ic.tensor\_contraction\_v1} are described in
\S\ref{sec:pt-tensor}; that arm has not been run.

\section{Black-box balance}\label{sec:tn-balance}

The lead note of 31 October
(\art{research/ic_prime_eth_leads_20261031/LEADS.md} \S\S2, 4) puts the
route in the cost model of the rest of this treatise. Let $s=b/m$ be the
window, $\abs F=2^s$ and $2^{es}$ the per-attempt decomposition cost. Three
statements follow, with Lean receipts in \texttt{lean/IcTensorRoute.lean}
(standard axioms, no \texttt{sorry}) that cover the cost-model arithmetic
only.
\begin{itemize}
\item \emph{Ladder rule.} The relation phase $m!\,p^{(1+e)/m}$ is below
  $p^{1/2}$ iff $e<m/2-1$, the gate of \S\ref{sec:pg-gate}.
\item \emph{Solver-saving floor.} An encoding has at least $b$ interface
  bits; a Boolean solver costing $2^{n_{\rm eff}(1-\delta)}$ meets the
  per-attempt budget only if $\delta\ge1/2+1/m$. Tensor contraction under
  D.1 gives $\delta=1/\sqrt{\log s}$, far below.
\item \emph{Black-box balance.} Letting $m$ grow trades the $1/m$ of the
  budget against the $m$ of the edge count; under the cost model
  $E(m)=b/m+cmb/g$ with $g=\sqrt{\log_2b}$, the optimum is
  $m\approx(\log b)^{1/4}$, and a route below rho needs $g>16c$, that is
  $\log_2b>256c^2$, with $c$ the unspecified constant in D.1.
\end{itemize}
The note adds a concrete size: at $b=256$ the balance gives $m=2$, and the
linear algebra alone costs $r^{3/2}\approx2^{384}$ against rho's $2^{128}$;
beating rho on linear algebra needs $m>6$, which the balance reaches only
for $b$ above about $2^{1296}$ (derived). The remaining levers it names are an
improvement of the contractor's exponent on these particular networks, the
open problem \art{problems/correlated-packet-common-kernel.md}, relation
networks with $o(mb)$ edges, and structure that lets a solver skip a
constant fraction of the interface.

\section{Twelve prime-field leads and their refutations}\label{sec:tn-leads}

The same note proposed twelve measurable leads, L01--L12, each with a
decision rule. They were evaluated on 31 October and reviewed by two
adversarial skeptics per lead; the skeptic digest is authoritative on
standing, and the synthesis reconciles the two
(\art{research/ic_prime_eth_leads_20261031/SYNTHESIS.md};
\art{research/ic_prime_eth_leads_20261031/leads/SKEPTIC_DIGEST.md}). The raw lead
directories are on an external volume, and nothing in the synthesis was
re-run. This section is the full account: it gives the standing per lead,
what the leads measured and why six were refuted as stated;
\S\ref{sec:pa-sat} summarises it for the prime-field chapter. The verdict of the synthesis is quoted: ``Nothing
moved the exponent question for prime-field ECDLP index calculus. Every
result is a constant, a known negative, or an unfinished harness.''

The errata of 2 November list the standing after review
(\art{research/ic_prime_eth_leads_20261031/LEADS.md} \S9): L03, L07 and L11
narrow, constant only; L04 and L09 contested and partial; \emph{L01, L05,
L06, L08, L10 and L12 refuted as stated}; L02 not run. For L05 the skeptic
digest is more precise: what is refuted is the lead's claim to be
\emph{closed} (``refuted as closed'', both skeptics); the harness and its
numbers stand.

\begin{table}[t]
\centering\small
\begin{tabularx}{\linewidth}{@{}>{\raggedright\arraybackslash}p{0.17\linewidth}XX@{}}
\toprule
Lead & As proposed or first read & After skeptic review \\
\midrule
L01 char-2 gadget & a Sch\"onhage-type gadget for $\FF_{2^n}$; evaluator verdict ``undecided'' because $\ell=7$ was censored & three finished $\ell=7$ runs give slope about $2.58$ bits of conflicts per window bit (reopen bar: CI below $2.5$; ladder bar $0.5$); the early ratio fall fits a constant padding floor of about 11k conflicts; CPU $3\times$ worse at $\ell=7$; only 4 of 11 products use the gadget \\
L05 policy pilot & closed after a pilot of an LLM policy-search harness & refuted as closed: renaming is a null experiment by construction in kissat; the ladder held $b=8$ and varied $\ell$, not $e$ against $s$; 3 curves; the hooked CaDiCaL engine was never built; label ``harness built, renaming null recorded, policy question untested'' \\
L06 contraction width & closed: tree-decomposition width loses to enumeration & measured on a Python structural model, not a graph dump; dense gate-graph treewidth is the wrong proxy for D.4; the transform-word cut was effectively not run; surviving claim: plain tree-decomposition DP on the school CNF loses at $b\le16$ (lower bound 15--51 against 2--10.7) \\
L08 Koblitz transfer & char-2 gadget no worse than the integer one, margin $3.42\times$ & that compared a freed binary gadget with an unfreed integer one; the skeptic's freed 133-bit integer multiplier has 14{,}618 edges instead of 37{,}030, so the margin is about $1.4\times$; hand bound $80d+228\sqrt d$ fails at $d=46$; supported bound about $175d$ for $d\le5000$ \\
L10 Lean filings & ``done; file as known results'' & seven gaps between the lead text and the Lean statements, all confirmed: the cost model is an unproved hypothesis, \texttt{no\_crossover} needs $c\ge1$, the interface bound is an injectivity count; not fileable as drafted \\
L12 clause retention & incremental solving amortises learned clauses, $1/R=40\times$ at $m=2$, $b=12$ & the retained solver fills an $\abs F^{m-1}$ table whose set-up is never charged against the explicit enumeration reference; fresh baseline censored at $m=3$; target lists 48\% duplicated; class retained set-up/advice, closed \\
\bottomrule
\end{tabularx}
\caption[]{The six leads refuted as stated. Producer-reported, from the
skeptic reconciliation; no lead has a validation record. Source:
\art{research/ic_prime_eth_leads_20261031/SYNTHESIS.md} \S3;
\art{research/ic_prime_eth_leads_20261031/LEADS.md} \S9.}
\label{tab:tn-refuted}
\end{table}

Table~\ref{tab:tn-refuted} gives the six refutations. The other five run
leads carry numbers that are worth keeping as harness constants.
\begin{itemize}
\item \emph{L03, sparse-parity bases.} Every hashed base (Bochum's affine
  hash, dense, sparse weight 3, 5, 7, biregular, window) matches a
  size-matched null to a paired ratio $1.000\pm0.0003$ at $p\approx2^{14}$ to
  $2^{22}$, $m=2,3$. Weight 3 loses about 10\% at $m=3$ from empty or
  heavy-tailed bases. This restates the conservation of
  Chapter~\ref{ch:cover}: a size-matched null makes pairwise-independent
  hashes match by construction.
\item \emph{L04, backdoor census.} The offset $(m-1)s-B^*$ saturates at
  $2\log_2(\mathrm{cap})-1$ (1, 3, 5 for caps 2, 4, 8) at $m=3$, $s=4..7$: the
  enumeration has moved inside the propagator. The supportable sentence is
  ``no sub-$(m-1)s$ strong backdoor under cap-bounded window-branching
  propagation at $s\le7$''.
\item \emph{L07, second moment.} Hashed bases are $1.000\pm0.005$ of the
  null; a shared negation-closed $x$-window gives $0.713/0.833/0.916$ at
  $m=4$, the symmetric-base mechanism already known.
\item \emph{L09, census tool.} Over 34 encodings the margin to D.1
  admission is $-16.11$ to $-16.33$ bits of $\log_2\lambda$; with D.1's
  constant set to 1, $s/\sqrt{\log_2s}$ is $37\times$ to $41{,}000\times$ the
  enumeration exponent.
\item \emph{L11, multiplier encodings.} Propagation-complete against school
  multiplication: paired ratio $0.78$ $[0.67,0.92]$ at $b=9$ only, $1.15\times$
  worse at $b=6$, untested at $b\ge12$.
\end{itemize}
The synthesis also records feasibility: $m=3$ prime SAT is censored above
$b\approx12$ in every lead that tried it, so only the $b=6$ and $b=9$ diagonal
cells are uncensored, and its closing recommendation is that the SAT and
propagation lane for prime fields ``should be closed for exponent
purposes'', with one item untested (a CaDiCaL external propagator with
$\FF_p$ Gaussian elimination).

\section{Corrected sentences of the lead note}\label{sec:tn-corrections}

The lead note is left unchanged so that its original claims stay auditable;
its \S9 lists eight corrections and says that where the text conflicts,
the errata win. Three of the corrected sentences concern the route itself.
\begin{itemize}
\item \emph{``The Bochum route can never beat rho on any curve anyone will
  deploy, whatever D.1's constant is''} is wrong as stated. Lean proves
  $g>16c$; the figure of more than $2^{256}$ bits needs $c\ge1$, an outside
  assumption on an unspecified $O(\cdot)$; the general threshold is
  $\log_2b>256c^2$.
\item \emph{``By the paper's own gadget counts the network is at least
  $\sim10^5b$ edges per point addition''} is $9$--$30\times$ high. Measured:
  3{,}364 edges per slot-bit for the freed binary network, $1.17\cdot10^4$ for
  the unfreed prototype, and 11.9k edges per bit for the $m=2$ prime network
  at $b=131$. The exponent argument is unaffected.
\item \emph{The clustering loophole} (merging vertices to reduce $s$ below
  the admission threshold) is closed: no row is D.1-admitted, hidden
  interface bits cost $2^h$ entry evaluations, $s_{\rm eff}\ge(h/A)^{1.11}$,
  and the $b<2^{36}$ closure is specific to $m=3$ ($b<2^{16}$ at $m=4$); $m\ge4$
  is closed by \texttt{no\_crossover}, not by the floor. The derivation, with
  heuristic steps labelled, is in \texttt{notes/clustering\_closure.md}.
\end{itemize}
The other five corrections are the L10 Lean overclaims, path errors, the
ladder feasibility limit, the L08 hand bound and the L08 margin; the last is
in Table~\ref{tab:tn-refuted}. The atlas carries the same corrections
(Table~\ref{tab:corr-exotic}). The skeptics refuted all nine follow-up proposals as proposed, in
one sentence each, and kept six cheap residues of them as harness constants;
none touches an exponent.

\section{Finite size}\label{sec:tn-finite}

The finite-size review of 30 October
(\art{research/hashed_ecdl_finite_size_20261030/README.md}) asked whether
the hashed relation language can be useful below the recursion's onset. It
argues that the onset is a guard in one algorithm, not a lower bound for
solving these particular networks, and proposes five search strategies
(listed in \S\ref{sec:pt-tensor}); nothing was implemented or measured
under it. Its fourth route, replacing the global recursion schedule with
checked local plans, is the closest to lowering the same recursion's onset
and, in the review's words, ``has the largest proof gap''. The state
assessment of the same day
(\art{research/hashed_ecdl_leads_20261030/README.md}) put the priorities as
an independent validation of the general contraction route and a
complete-cost evaluation of the finite algorithms already delivered; it
records the finite hashed track as implemented and development-tested
(selectors, last-slot derivation, width and cut-set plans, partial
relations, binary components) with ``no admitted matched full-cost panel or
evidence of an exponent improvement''. The elimination campaign of
1 November and the rank-census design are in \S\ref{sec:pt-tensor}.

\section{What the route has and has not shown}\label{sec:tn-reading}

The route is the only claimed sub-exponential algorithm for prime-field
ECDLP that the programme examined in detail, and it is external. What the
programme established is:
\begin{itemize}
\item an independent native implementation of the complete algorithm,
  checked on small curves, with an exact contraction and certificate
  replay;
\item that the literal constants select brute force on every network that
  can be written down, and that both published challenges exceed the exact
  budget without a relation;
\item that under the lead note's cost model black-box use of D.1 needs
  $\log_2b>256c^2$ to cross rho, with Lean receipts for the arithmetic;
\item twelve leads whose measured outcomes are constants or known
  negatives, six of them refuted as first stated.
\end{itemize}
What it has not established is any counted cost, any exponent, or the
correctness of D.1. The route does not move a row of any table in this
treatise. It remains open in the atlas with two next steps: the independent
audit of correctness and complexity, and the frozen 84-case comparison
panel, which is blocked on checkpoint transport and native execution access.
If D.1 were withdrawn, the route's asymptotic claim would go with it and the
implementation would remain an exact enumerator with a tensor interface.

% ---------------------------------------------------------------------------

\chapter{Learning, pseudorandomness and curve models}\label{ch:learning}

This chapter collects three kinds of unconventional work that share one
methodological lesson: an instrument that sees nothing excludes only what
it could have seen. The first is learning inside the decomposition oracle,
where a preregistered study measured the most a learned component could gain
and found the gain unreachable. The second is a set of tests of whether
the factor base looks random in discrete-logarithm order, and a calibration
that showed which known-weak curves those tests can and cannot see. The
third is the remainder of the lifts, transports and curve models not
already in Chapter~\ref{ch:obstructions}.

\section{The ceiling for a learned decomposition oracle}\label{sec:ln-ceiling}

Chapter~\ref{sec:binary} records that SAT, SMT and WDSat are enumerators on
binary point decomposition. A study of 27 October
(\art{research/pdp_ml_sat_ceiling_20261027/README.md}; protocol frozen at
\texttt{8a85584fd} with one amendment before any result) asked the question
a learned component would have to answer before anyone builds a graph
network or an RL branching policy: how much could any learned component
gain inside a SAT point-decomposition oracle, and does any cheap signal
reach it? The setting is the symmetrised $\Sem4$ at $m=3$, a native
CDCL$+$XOR solver, the curves of the Trimoska corpus~\cite{trimoska2020} and
uniform targets, in two families: $n=15$, $\ell=5$ (3{,}000 targets) and $n=17$,
$\ell=6$ (1{,}200 targets). Cost is in solver conflicts. Three levers were
tested: triage (finish only the solves that score high), a fixed conflict
cutoff, and branching order.

\begin{table}[t]
\centering\small
\begin{tabularx}{\linewidth}{@{}Xll@{}}
\toprule
Quantity & $n15\ell5$ & $n17\ell6$ \\
\midrule
UNSAT conflicts, median [IQR] & 5781 [5738, 5819] & 45649 [45539, 45781] \\
conflicts per sorted factor-base triple & 0.963 & 0.998 \\
perfect-predictor ceiling, relations per conflict & 5.21--9.01$\times$ & 5.75--6.69$\times$ \\
test AUC, all feature groups & 0.485--0.556 & 0.461--0.513 \\
permuted-label AUC & 0.478--0.540 & 0.449--0.517 \\
positive-control AUC (label flipped w.p.\ 0.25) & 0.747, 0.785 & 0.717, 0.725 \\
fixed-cutoff test gain & 0.93, 1.05 & 0.91, 0.98 \\
learned-activity branching order, conflicts / default & $\ge10.03$ & $\ge10.00$ \\
\bottomrule
\end{tabularx}
\caption[]{The ML-aided SAT ceiling for point decomposition (symmetrised
$\Sem4$, $m=3$, binary toy families). Ranges are over the two splits and
four feature groups. Producer-reported; toy stage diagnostic, no
validation record. Source:
\art{research/pdp_ml_sat_ceiling_20261027/README.md}.}
\label{tab:ln-ceiling}
\end{table}

Table~\ref{tab:ln-ceiling} gives the numbers. All 4{,}200 labels agree with
exhaustive sorted-triple search. A refutation costs about one conflict per
sorted factor-base triple, and every refutation lands within 4.1\% of its
family median, so the solver is an enumerator. The prize is real: a perfect
UNSAT predictor, charged its telemetry conflicts, would raise relations per
conflict by $5.2$--$9.0\times$. Nothing reaches it. Test AUC sits within the
noise of $0.5$ for every feature group (bits of $x_r$, their pairwise
products, solver telemetry after 64 or 512 conflicts), on both families and
both splits; the positive control at $0.72$--$0.79$ shows that the learner
works. The one interval excluding $0.5$, telemetry at 512 conflicts on
$n15\ell5$ at $0.556$ $[0.518,0.599]$, includes the ``decided within 512
conflicts'' flag and gains $0.99\times$. Train-fitted gains of up to $4.7\times$
fall to about $1\times$ held out, and telemetry trained at $n=15$ scores $0.50$ and
$0.52$ at $n=17$. A cutoff gains nothing because a witness found early is no
more likely than one found late. Default VSIDS is at the floor: an order
learned from VSIDS activity on training refutations is at least $10\times$
worse on every held-out target, because the variables that take part in
conflicts are not the ones worth deciding on.

The study closed under its preregistered rule (no lever reached $1.5\times$
with a lower bound above $1.2\times$ on both families; the largest held-out
lower bound was $0.92$). Its prediction P4, that every branching arm would
stay within $0.7$--$1.4\times$ of the default, failed. The reading for index
calculus is the product law the README states: even a perfect classifier is
capped by relations $\times$ targets $\times$ oracle $\approx m\cdot2^n$, so it removes
a constant from the oracle (measured here at $5$--$9\times$) and cannot change
the exponent. The run also found and fixed a soundness bug in the native
solver: resuming after a budget stop answered wrongly in both directions;
an audit found no published artefact that re-solved without a reset.

\section{The ML dataset campaign}\label{sec:ln-dataset}

From 28 October to 1 November the programme built a native
\texttt{crypto dataset} generator and custody pipeline for a corpus on
which to train models of ECDLP and point-decomposition behaviour: Koblitz,
ordinary binary and prime curves with paired twists, arities $m=2$ to 15, bit
sizes 23--31 and 47--61, multiple generators and keys, and every supported
native method (\art{docs/ml-dataset/DESIGN.md}; more than forty
\texttt{research/ml\_dataset\_*} directories). The design's planning
arithmetic is 2{,}016 parameter cells and, at 16 curve seeds, 8 generators
and 64 keys, 16{,}515{,}072 tasks per backend configuration; the stratified
design keeps 165{,}150{,}720 method cells at full width. These are planned
cells, not data.

\paragraph{Scope.} The campaign inventory of 1 November
(\art{research/ml_dataset_campaign_inventory_20261101/README.md}) records
124 selected terminal method cells in seven cohorts, under different
protocols, seeds, methods and caps, and states that the count ``is a
coverage inventory, not a pooled estimate or a training set''. Its labels are
mostly method completion, bounded candidate counts and censored statuses.
Examples: Koblitz $n=23$, $m=3$--15, 26 cells with 4 finite-domain negatives
and 22 timeouts; Koblitz $n=47$, $m=4$, 16 verified method artifacts;
$p=2^{61}-1$, $m=2$, 2 verified enumerations and 2 root timeouts; $n=26$--31,
$m=2$, 36 cells, 35 known-answer-replayed. A separate 252-cell compatibility
matrix is development-only. The one model study
(\art{research/ecdlp_dataset_models_20261102/REPORT.typ}) is a development
pilot that by its own protocol may not select a winner or estimate
generalisation: 960 cells, 819 training rows, and an arm that nearly fixes the
label (\texttt{enumerate} completes on all 240 rows and recovers the scalar
on none; \texttt{bsgs} solves 237 of 240). Only the double-large-prime arm
shows signal, a Brier score of $0.050$ against $0.192$ for the majority
baseline, which the permutation control removes ($0.168$). It ran in the
session container, not through the experiment route.

\paragraph{Relevance to index calculus.} The campaign has no model result,
and it bears on index calculus only through the ceiling of
\S\ref{sec:ln-ceiling}: a learned predictor of decomposition success can at
best remove the constant measured there. What it does contribute is a label
discipline that the rest of this treatise also follows
(\art{docs/ml-dataset/DESIGN.md}, ``Scientific labels''):
\begin{itemize}
\item target creation by $Q=[k]G$ is not an independently solved ECDLP, and a
  verified decomposition is not an ECDLP solve;
\item the first degree drop in one algorithm's trace is not the first fall
  degree, and the maximum degree visited is a solving-degree observation, not
  the last fall degree (Chapter~\ref{ch:falldegree});
\item censored outcomes (timeouts, budget exhaustion, unsupported cells) are
  kept, not dropped;
\item splits group all keys and generators on one curve, both twists,
  isomorphic and isogenous models and Frobenius-equivalent tasks into one
  component; random row splits are inadmissible.
\end{itemize}

\section{Other learned components}\label{sec:ln-components}

Four smaller studies tested learned components directly; none improved on
its incumbent.
\begin{itemize}
\item \emph{Learned WDSat branching}
  (\art{cryptanalysis:experiments/ideas_learning_20261018/REPORT.md}, item
  8). 88 solver calls with upstream WDSat \texttt{61c6ff3f}: on eight held-out
  targets the paper default with symmetry took 3{,}092 conflicts, and the
  block, interleaved, random and learned static orders without symmetry took
  16{,}097, 16{,}785, 16{,}676 and 16{,}097. The learned policy selected the
  ordinary ascending order, so it \emph{is} the block arm. The
  $16{,}097$-against-$3{,}092$ gap is the symmetry switch, which the controls
  disable because the upstream symmetry code interprets fixed positions in
  the order; it is not a learned-policy effect
  (Table~\ref{tab:corr-exotic}). The atlas still quotes the comparison as
  ``the only learned-branching negative''; read correctly, it shows that the
  learned order did not improve the ordinary order.
\item \emph{Learned pivot selection and program search}
  (\art{research/ideas_execution_20261018/EXECUTION.md}, items 9--10). A
  native learned F4 pivot selector returned the constant minimum-weight
  rule. A restricted search over group-law query kernels re-selected the
  incumbent pair-hash kernel; both took 820 counted units on the held-out
  curve.
\item \emph{Learned prefilter} (item 11). On 7{,}254 exact subgroup
  targets in 27 cells, the one-class prefilter rejected genuine
  decompositions and was disabled.
\item \emph{Degree surrogate} (item 12). Fitted to 1{,}718 archived
  degree-profile cells, the curve-held-out surrogate has degree MAE $0.10041$,
  but with three calibration curves the honest 90\% group-conformal interval
  is unbounded, so it cannot price a future run.
\end{itemize}
A reinforcement-learning pair selector for Buchberger's algorithm
(\art{research/gbrl/PAPER.md}, undated, a paper-style self-report with no
frozen protocol or ecbench record) reports $-5.5\%$ S-polynomial steps on
cyclic-6 and, on held-out Semaev $m=2$ systems, $-3.9\%$ steps with $+5.4\%$
arithmetic operations; at $m=3$ it fails to generalise to a factor base of
size 6 (496 steps against sugar's 459). It works in Buchberger only, not F4.
The L05 pilot of Chapter~\ref{ch:tensor} closed LLM policy search on variable
renaming. The atlas's next step for the whole learning lead is a
learnability control: predict a log bit from coordinates with a planted
homomorphic leak as positive control and a random-relabelling null, at $N$
up to $2^{24}$.

\section{Pseudorandomness tests}\label{sec:ln-instruments}

The question behind this line is whether a natural factor base, ordered by
discrete logarithm, looks like a random set. Any detectable structure would
be a lever: a relation yield above the counting heuristic, or a statistic
that leaks log bits. The linear-phase test, H-PSEUDO, is told in
Chapter~\ref{ch:summer}, with its 4 November correction: the small-$x$ base
is negation-closed and the original null was not, and against a
negation-closed null $R$ is consistent with 1 (Table~\ref{tab:corr-xref}).
Three further instruments followed, each under a protocol frozen before
data.
\begin{itemize}
\item \emph{Quadratic phases} (QUADPHASE,
  \art{crypto-autoresearcher:docs/reports/quadphase-20261104/README.md}). The
  statistic is the maximum over $(\alpha,\beta)$ of
  $\abs{\sum_{l\in S}e((\alpha l^2+\beta l)/N)}/\sqrt{\abs S}$ for the logs $S$ of the
  first 512 points in $(x,y)$ order, against 40 negation-closed random sets
  per curve. Verdict NULL: 9 of 9 decisive curves at $N\approx2^{12}$--$2^{14}$
  inside the null's band ($u$ from $0.125$ to $0.85$); planted control
  $3.8$--$5.2\times$ the null maximum at the planted $\alpha_0$; A/A difference at
  most $2.3\%$ pooled. An independent validation recomputed all 385 rows with
  an FFT-free implementation to within $7\cdot10^{-15}$. The validator's probe
  says the test detects coherent single-phase leakage of about $0.3$ of the
  set or more; weaker or multi-peak structure is untested. One descriptive
  curve at $2^{16}$ sat above all 5 of its null draws; a rerun against 40 draws
  (\art{crypto-autoresearcher:docs/reports/quadphase-2p16-20261104/README.md})
  put it at the median, $u=0.50$ (null mean $4.750$, s.d.\ $0.229$, maximum
  $5.581$).
\item \emph{Digit-level tensor structure} (HOLANT-STEPB-QTT,
  \art{crypto-autoresearcher:docs/reports/holant-step-b-qtt-20261104/README.md}).
  The holant lead asked whether a basis change on the $\FF_p$ bond alphabet
  makes the decomposition-count network cheap. Flattening and CP ranks are
  invariant under invertible bond-wise maps (checked at $p=61$), so the test
  was reformulated as the quantised tensor-train rank of the transfer matrix
  $M$ under binary digit splitting. NULL: from $L=8$ the natural
  $x$-labelling and the random relabellings both sit at the rank ceiling
  $2^L$ (1024 at $p=1021$), while the discrete-log labelling, the positive
  control, has rank 109--121. A low score on one anomalous curve at $p=61$
  was followed up on 63 anomalous and 63 ordinary curves at $37\le p\le127$
  (\art{crypto-autoresearcher:docs/reports/anomalous-qtt-20261104/README.md}):
  NULL, permutation $p=0.831$, the anomalous curves slightly above the
  ordinary ones.
\item \emph{Coordinate statistics.} The census that found every structured
  additive statistic $\mathrm{PGL}_2$-unstable and every stable multiplicative
  coset statistic at the random level (\texttt{EV-ECDLP-81f4d4}) was lowered
  from strong to preliminary by \texttt{CORR-20261026-c94668}
  (Table~\ref{tab:corr-structure}). The record's next step is the one
  natural family the census never measured, digit characters.
\end{itemize}
Together with H-PSEUDO, these say that at toy size the small-$x$ factor base
in log order looks like a random negation-closed set under linear and
coherent quadratic phases and under digit-level tensor rank. They say it
only for structure of each statistic's type, and the next section shows how
narrow that is.

\section{Calibration against known-weak curves}\label{sec:ln-weak}

The weak-class calibration of 4 November
(\art{crypto-autoresearcher:docs/reports/weakclass-20261104/README.md},
WEAKCLASS-20261104) asked whether these instruments see curves whose
discrete logarithm is \emph{known} to be easy: anomalous curves ($n=p$,
Smart's attack), small embedding degree ($\mathrm{ord}_n(p)=3$--6, MOV), and
$j=0$ curves with a factor base of whole $\zeta$-orbits, each with a paired
ordinary curve over the same prime, with $p$ of 10, 12 and 14 bits and 3
weak curves per class per size. A cell
detects only if all 3 class curves fall beyond the null's 5--95\% band in the
same direction at two sizes and the paired ordinary curves do not.

\begin{table}[t]
\centering\small
\begin{tabularx}{\linewidth}{@{}Xlll@{}}
\toprule
Instrument and run & anomalous & small $k$ & $j=0$, $\zeta$-stable base \\
\midrule
I1 character sum, WEAKCLASS & invalid ($\bar u=0.61$) & invalid ($0.41$) & invalid ($0.12$) \\
I1 rerun, calibrated A/A & blind ($0.46$) & blind ($0.44$) & blind ($0.28$) \\
I1, 36 curves per class & not detected ($0.483$, $p=0.86$) & not detected ($0.476$, $p=0.82$) & detected ($0.233$, $p<10^{-5}$) \\
I2 $m=2$ sumset yield & blind ($0.56$) & blind ($0.44$) & blind ($0.64$) \\
I3 QTT rank, $p\approx2^{10}$ & blind ($0.36$) & blind ($0.44$) & blind ($0.60$) \\
\bottomrule
\end{tabularx}
\caption[]{Which instrument sees which known-weak class. Entries are the
verdict and mean percentile $\bar u$ within the instrument's own null. All
producer-reported; the third row's study names an independent validation
that replayed all 324 rows. Sources:
\art{crypto-autoresearcher:docs/reports/weakclass-20261104/README.md};
\art{crypto-autoresearcher:docs/reports/weakclass-i1-20261104/README.md};
\art{crypto-autoresearcher:docs/reports/paired-anomaly-20261104/README.md}.}
\label{tab:ln-weak}
\end{table}

Table~\ref{tab:ln-weak} gives the outcome. I2 and I3 are blind to all three
classes while their positive controls sit at $u=0.00$ on every curve, so the
instruments work and simply do not see Smart- or MOV-weak curves. I1 came out
\emph{invalid} in the first run: its preregistered A/A gate required even and
odd null means to agree within 5\% on every curve, while with 20$+$20 draws the
sampling error alone is about 2\%, and 5 of 54 curves exceeded it. The study
attributes the fault to the gate, not the instrument
(Table~\ref{tab:corr-structure}). The rerun
(\art{crypto-autoresearcher:docs/reports/weakclass-i1-20261104/README.md})
used 200 null draws per curve in two streams compared by a KS test (7 curves
rejected against an allowance of 14) and found I1 valid and blind to all
three classes. It also found a post-hoc anomaly: the paired ordinary curves
sat low (mean $u=0.276$, KS $p=1.4\cdot10^{-4}$), which put the pairing baseline in
doubt. A preregistered replication
(\art{crypto-autoresearcher:docs/reports/paired-anomaly-20261104/README.md},
108 paired and 108 independent curves) did not replicate it (Mann--Whitney
$z=0.45$, $p=0.65$), and with 36 curves per class it gave the one detection
in the table: the $\zeta$-stable factor base of the $j=0$ class sits low at every
size (mean $u=0.233$, $z=-4.79$). The study explains why: a log set closed under
multiplication by the eigenvalue $\lambda$ has fewer independent frequencies,
so its character-sum maximum is lower. That reflects how the factor base was
built, not how easy the discrete logarithm is.

The consequence is a scope note on every null above. With 36 curves per class
at $p\le2^{14}$, I1 shows no detected shift for anomalous or small-embedding-degree
curves, with a minimum detectable shift of about $0.13$ in $u$ at 50\% power;
I2 and I3 are blind at 3 curves per class, I3 only at $p\approx2^{10}$ where it is
rank-saturated. So the instruments that found the small-$x$ base random in log
order would also have found it random on curves where the discrete logarithm
is polynomial-time or sub-exponential. A NULL from them excludes structure of
the statistic's own type and says nothing about whether a curve's discrete
logarithm is easy. The methodology lead's next step, injecting planted
signals into the real pipeline to measure each screen's false-negative rate,
has not run.

\section{Lifts, transports and curve models}\label{sec:ln-models}

Chapter~\ref{ch:obstructions} records the transport, lifting and twist
obstructions, and its xedni section is now read as conditional; the atlas
reopened the global-lift leads accordingly, with short Mordell--Weil
relations found on function-field surfaces at $p\le31$ but no decomposition
step. The remaining lift and model records are short.
\begin{itemize}
\item \emph{$p$-adic jets and canonical lifts.} For $\gcd(n,p)=1$ every additive
  order-$\ell$ image of a jet or formal-group target is zero (an independently
  verified scoped gate). The digit identity $d(m\hat S)=m\,d(\hat S)$ holds exactly,
  1360 of 1360, but only on the hard global lift, whose cost is the lift itself;
  the computable Teichm\"uller contrast reads as chance, and a valuation profile
  collapses to the $\pm$ fibre indicator (\texttt{KN-FIND-85845f},
  \texttt{KN-FIND-e8cfdc}, \texttt{KN-FIND-1ba0fe}). The September statistic is in
  Chapter~\ref{ch:summer}.
\item \emph{Elliptic nets.} Somos identities hold for every $k$, so net
  relations encode the group law on one fibre and not $k$. On toy primes
  $p\in\{101,431,1601\}$ the Somos sieve's charged $k$-recovery exponent is
  $1.97$--$2.80$ against BSGS's $1.17$--$1.32$ at identical accounting, and the
  sieve-to-BSGS operation ratio grows from about $24\times$ to $135\times$
  (\texttt{EV-NET-001}, preliminary).
\item \emph{Curve models.} Model-native solving degree is the same for
  Weierstrass Semaev and twisted Edwards systems, $d_{\rm reg}=2$ at $m=2$ and 3,
  with Edwards a constant factor larger (\path{EV-REP-001},
  \path{EV-REP-002}, replicated, toy). The EcGFp5 representation audit and
  its correction \texttt{CORR-20261102-7e3a91} are in \S\ref{sec:pt-gfp5}.
\item \emph{Split Jacobians.} Embedding $E$ as a Poincar\'e factor of a genus-2
  Jacobian cannot beat rho on the elliptic factor because the Hasse--Weil order
  inflation of the ambient Jacobian outweighs the genus-2 index-calculus saving
  (\texttt{KN-FIND-61347e}); the projected-smoothness counting gate finds no
  route with $\lambda,\mu\le0.45$.
\item \emph{Torsors and deck orbits.} For a generically Galois cover every
  rational deck invariant factors through the quotient, and a connected
  nontrivial cover has no rational section, so a target-compatible selector is
  a missing section (independently verified scoped gate).
\item \emph{Miller cocycles.} When $\gcd(N,p-1)=1$ the $N$-th root of a Miller
  function is a canonical primitive computable from the point; whether it
  exposes $k$ cheaply is an unreviewed proposal with no measurement.
\item \emph{Latt\`es dynamics.} Every forward iterate on the $x$-line is a group
  operation, so forward dynamics is generic, and orbit quotients give
  $\sqrt{\abs{\mathrm{Aut}}}$ for units and nothing for non-units; two accumulator
  closures are recorded (\texttt{EV-ECDLP-0c0f51}, \texttt{EV-ECDLP-a8af77}).
\item \emph{Lossy tracked objects.} An orbit argument shows that no lossy
  object propagates deterministically under the full translation action, a
  rigidity theorem covers single-point objects, and the one measurement, at
  $\FF_{17}$, is mixed (\texttt{EV-ECDLP-ce9b50}).
\end{itemize}
Each record is in its YAML file under
\art{crypto-autoresearcher:knowledge/leads/ecdlp/exotic/}; none has a cost
against rho.

\section{Out of scope}\label{sec:ln-scope}

Two lines of work in the same repositories are not classical ECDLP index
calculus and are excluded from this treatise's accounting.
\begin{itemize}
\item \emph{Quantum resources.} The reversible secp256k1 point-addition
  circuit for Shor's algorithm records a best score of 2{,}490{,}847{,}962
  (average Toffoli count 1{,}736{,}993 times 1{,}434 qubits), $-0.124\%$ against
  its head, valid on a 9{,}024-shot sample island only
  (\art{crypto-autoresearcher:knowledge/leads/ecdlp/exotic/LEAD-EXO-quantum-ecdlp-resources.yaml}).
  It is a resource estimate for a quantum attack, not a classical cost.
\item \emph{Eigen-S\'eta} is an isogeny-based public-key encryption scheme
  (\art{research/eigen_seta_assessment_20261017/REPORT.md} and related
  directories), not an ECDLP method; its corrections are listed in
  Table~\ref{tab:corr-exotic} for completeness. The neural distinguishers on
  Speck and AES (\art{research/neural_distinguishers_20261102/}) are symmetric
  cryptanalysis.
\end{itemize}

\section{What the three chapters settle}\label{sec:ln-reading}

None of the unconventional routes produced a counted cost of a discrete
logarithm, so none changes the distance to rho of Chapter~\ref{ch:distance}.
What they settle is narrower and mostly methodological.
\begin{itemize}
\item Generic and query bounds do not constrain index calculus, because
  they leave local work with coordinates free: the SGGM is vacuous on 122 of
  208 configurations, every coordinate oracle is outside Shoup's model, and
  the low-degree bound holds only for small multipliers
  (Chapter~\ref{ch:unconvmap}).
\item The tensor route rests on an unverified external theorem, is brute
  force at every realisable size, and found no relation on ECC2K-130 or
  ECCp-131; six of its twelve prime-field leads were refuted as stated and
  none moved an exponent (Chapter~\ref{ch:tensor}).
\item A perfect learned predictor would buy $5$--$9\times$ in a SAT
  decomposition oracle at toy size and nothing in the exponent; no feature
  tested reaches it (\S\ref{sec:ln-ceiling}).
\item The pseudorandomness instruments find the small-$x$ base random in log
  order at toy size, and they are blind to anomalous and small-embedding-degree
  curves; the only class they detect is a factor base built with a symmetry
  (\S\ref{sec:ln-weak}).
\end{itemize}
The corrections these chapters state are: the learned-WDSat comparison is
confounded with the symmetry switch; the I1 instrument was invalid in its
first run because of a miscalibrated A/A gate; the paired-curve anomaly did
not replicate; L01, L05, L06, L08, L10 and L12 were refuted as stated; three
sentences of the tensor lead note were corrected; the hashed-ECDL set-up hang
was misattributed; and the low-degree protocol's own reading of a toy-size
pass was narrowed by its validator. The open items are the independent audit
of the tensor backend and D.1, a counted cost for any tensor arm, a lower
bound that charges local work, the learnability control with a planted leak,
and the false-negative rate of the structure screens.
